\chapter{Conclusion and Future Work}
\label{chap:conclusion}

\section{Summary of Contributions}

This thesis has presented:

\begin{enumerate}
  \item \textbf{\gonzales{}}: a production path tracer, sole-authored,
    capable of rendering Disney's Moana island scene.
  \item \textbf{\borg{}}: the first open-source GPU backed by physical
    silicon tapeout, implementing TBDR with Morton-encoded tile buffers.
  \item \textbf{Ray tracing ISA extensions} enabling \borg{} to
    perform hardware-accelerated BVH traversal and triangle intersection.
  \item \textbf{Gonzales-on-Borg}: a demonstration of path tracing
    on fully open hardware -- reproducible, inspectable, extensible.
\end{enumerate}

We believe this work establishes the first genuinely open hardware
research platform for physically based rendering.

\section{Future Work}

\subsection{Short Term}

\begin{itemize}
  \item Full Moana island scene on \borg{} (requires tiled streaming).
  \item Wavefront path tracing to improve hardware utilisation.
  \item Vulkan 1.0 conformance layer on \borg{}.
\end{itemize}

\subsection{Medium Term}

\begin{itemize}
  \item Multi-tile (chiplet) \borg{} configuration via LiteX/PCIe.
  \item ReSTIR direct lighting in \borg{} hardware.
  \item ECP5 / Artix-7 ports for higher logic capacity.
\end{itemize}

\subsection{Long Term}

\begin{itemize}
  \item Full hardware BVH builder on \borg{}.
  \item Open-source competitor level comparison with NVIDIA Turing RT Core
    at equivalent process node.
  \item Use as pedagogical platform: a student can build a working GPU
    from this thesis in one semester.
\end{itemize}

\section{Closing Remark}

Eric Veach's 1997 dissertation established the mathematical foundations
of Monte Carlo light transport.
This thesis hopes to do for hardware what Veach did for algorithms:
provide a rigorous, reproducible, and freely available foundation
that others can build upon.
